\section*{ID9510: Concluding Synthesis: g(0)}
\paragraph{Metadata.}
PiID: 9086996033027 | RTEID: RTE:P36929.4420:T4.2997 | UUID: cbc464f6-9cbb-5329-abdc-e9d45c435d91

\paragraph{Canonical Statement.}
Finally, we consider synchronization phenomena – from the Kuramoto model of coupled oscillators to celestial mechanics (tidal locking). Here, dimensionless thresholds do appear, but 1.185 is not a universal constant either. The Kuramoto model’s classic critical coupling for identical oscillators with a unimodal frequency distribution is given by \$K\_c = \textbackslash{}frac\{2\}\{\textbackslash{}pi g(0)\}\$ math.leidenuniv.nl math.leidenuniv.nl , where \$g(0)\$ is the density of natural frequencies at zero. Depending on \$g(0)\$, \$K\_c\$ can take various values. For a common case (Lorentzian distribution), \$g(0)=\textbackslash{}frac\{1\}\{\textbackslash{}pi\}\$ and thus \$K\_c=2\$ math.leidenuniv.nl . Other distributions yield different \$K\_c\$ (no fixed “1.185” across all systems). We did find a specific scenario where a ratio \textasciitilde{}1.18 emerged: in a modified Kuramoto mean-field model, coexistence of standing-wave and partial-sync states persisted up to a frequency ratio \textasciitilde{}1.18, beyond which the system went fully coherent math.leidenuniv.nl . This is essentially an empirically observed bifurcation point for that model’s parameters, not a general law. It shows

\paragraph{Canonical Equation.}
\[
g(0)=\frac{1}{\pi}
\]

\paragraph{Dictionary.}
Finally, we consider synchronization phenomena – from the Kuramoto model of coupled oscillators to celestial mechanics (tidal locking) (g(0))

\paragraph{Encyclopedia.}
Concluding Synthesis: g(0) is an algebraic definition, written g(0)=(1)/(π). It is an algebraic definition expressing the quantity directly in terms of the variables on its right-hand side. It is built from π, defining g(0). Within the Trawin Zero Point registry it serves as one element of the framework's mathematical narrative.
